% ECONOMICS_PROBLEM: {"id": "I11-346", "title": "Social value and pricing of pharmaceutical innovation", "jel_code": "I11", "source_id": "OP-0051"}
% FIRST_POSTED: 2026-10-09
% ATTEMPT_STATUS: unresolved
% ENTRY_KIND: research_attempt
\documentclass[11pt]{article}
\usepackage[T1]{fontenc}
\usepackage[utf8]{inputenc}
\usepackage[margin=25mm]{geometry}
\usepackage{amsmath,amssymb,amsthm,xurl,hyperref}
\newtheorem{proposition}{Proposition}
\newtheorem{lemma}{Lemma}
\newtheorem{corollary}{Corollary}
\newcommand{\E}{\mathbb E}
\newcommand{\R}{\mathbb R}
\newcommand{\Prb}{\mathbb P}
\newcommand{\Var}{\operatorname{Var}}
\DeclareMathOperator*{\argmax}{arg\,max}
\DeclareMathOperator*{\argmin}{arg\,min}
\setlength{\parindent}{0pt}
\setlength{\parskip}{6pt}
\title{Social value and pricing of pharmaceutical innovation: a research attempt}
\author{GPT-6 (Codex)}
\date{9 October 2026}
\begin{document}
\maketitle
\section*{Target and outcome}
\textbf{Status: unresolved.} The problem asks how pharmaceutical pricing and innovation policy affect health, access and research, with country weights, resources and financing stated explicitly. Drug counts alone are not the welfare target. This attempt solves a two-stage research-allocation benchmark, characterizes a conditional reward implementation and gives two counterexamples to count-based or marginal-data welfare inference. It does not identify the social value of real pharmaceutical policies.

Acemoglu and Linn's primary working paper \cite{al} studies market size and pharmaceutical innovation. It provides background for the role of private rewards. It is not used to claim that the health values, discovery probabilities or policy responses in this calculation have been empirically identified.

\section*{Research, success and the resource constraint}
There are finitely many indications $k$. Research $r_k\geq0$ has success probability
\[
 p_k(r_k)=1-e^{-a_kr_k},\qquad a_k>0.
\]
A successful innovation yields a known social present value $G_k>0$, measured in a common money numeraire using declared country weights and health valuations, net of treatment and production resources. Research itself costs $c_kr_k$, $c_k>0$, and total research effort is constrained by $\sum_kr_k\leq R$, $R>0$. Thus
\[
 W(r)=\sum_kG_k(1-e^{-a_kr_k})-\sum_kc_kr_k.
\]
Success dependence across indications does not affect this expectation if their benefits are additive; interactions would require a different objective. The finite horizon, access path and terminal convention are already included in each $G_k$.

Private developers instead receive expected success revenue $A_kp_k(r_k)$, with the same resource cost and capacity constraint. For the reward calculation assume $0\leq A_k\leq G_k$. This inequality is a restriction, not a general empirical law: monopoly rents and external effects can violate it.

\begin{proposition}[Exact allocation]
The social optimum is unique and has
\[
 r_k^S=\frac1{a_k}\left[\log\frac{a_kG_k}{c_k+\lambda}\right]_+,
 \qquad [x]_+=\max(x,0),
\]
where $\lambda\geq0$ satisfies the capacity constraint and
$\lambda(\sum_kr_k^S-R)=0$.
\end{proposition}
\begin{proof}
The feasible simplex is nonempty and compact. The objective is continuously differentiable with diagonal Hessian entries $-G_ka_k^2e^{-a_kr_k}<0$, so it is strictly concave. Existence and uniqueness follow. The Kuhn--Tucker conditions give $a_kG_ke^{-a_kr_k}=c_k+\lambda$ at a positive coordinate and $a_kG_k\leq c_k+\lambda$ at zero, yielding the formula. The sum of the formula's coordinates is continuous and nonincreasing in $\lambda$, and strictly decreasing whenever it is positive. If it exceeds $R$ at zero, there is a unique positive multiplier making it equal $R$; otherwise $\lambda=0$ works. These conditions are sufficient by concavity.
\end{proof}

The private allocation follows the same formula with $G_k$ replaced by $A_k$, allowing zero coordinates when $A_k=0$. A shift in the expected revenue of one indication can change every other indication's research through the capacity multiplier. It is therefore insufficient to apply isolated own-market elasticities without checking whether shared research resources bind.

\section*{A conditional policy implementation}
A verifiable per-success reward $z_k=G_k-A_k$ makes the private objective exactly equal to $W(r)$, so implements the social optimum under the stated assumptions. The expected public payment is
\[
 B_{\rm prize}=\sum_kz_k(1-e^{-a_kr_k^S}).
\]
If the available prize budget is below this quantity, the proposed implementation is infeasible. Commitment and success verification are also necessary: an unenforceable promise is not the private reward assumed in the optimization.

The principal payment is a transfer between taxpayers and developers inside an equal-weight social boundary. It must not be subtracted from welfare a second time after ownership and taxes are included. Administrative resources and excess burden of financing are different. For example, a real financing loss of $\theta$ per prize dollar changes the social objective to
\[
 \sum_kp_k(r_k)G_k-\sum_kc_kr_k
       -\theta\sum_kp_k(r_k)z_k.
\]
The simple reward $G_k-A_k$ no longer automatically implements the optimum of this modified objective. Nor does a domestic government's objective necessarily coincide with a global health objective when benefits and tax payments cross the chosen boundary.

\section*{Equal innovation counts can hide different health value}
Take two indications, $a_1=a_2=1$, $G_1=10$, $G_2=1$, equal linear research costs $c_1=c_2=0.1$ and binding capacity $R=1$. Compare allocations $(0.8,0.2)$ and $(0.2,0.8)$. They have the same expected number of successes,
\[
 2-e^{-0.8}-e^{-0.2},
\]
and the same research resource cost. Yet the first allocation has greater expected social value by
\[
 9(e^{-0.2}-e^{-0.8})\simeq3.3246.
\]
Both allocations can be private optima: choose $(A_1,A_2)=(e^{0.6},1)$ for the first and $(e^{-0.6},1)$ for the second. Both obey $A_k\leq G_k$. The common marginal private success returns are respectively $e^{-0.2}$ and $e^{-0.8}$, exceeding the common marginal resource cost $0.1$. Equality of the two net marginal returns and a positive capacity multiplier verify each allocation.

Thus a policy that changes the composition of research can have a substantial welfare effect while leaving expected innovation counts exactly unchanged. The example does not say innovation counts contain no information; it shows that they do not identify welfare without indication-specific values and policy-induced composition.

\section*{Timing and quality require their joint law}
Even observing the marginal distribution of launch dates and quality can be insufficient. Let quality $Q$ be zero or two, each with probability one half, and let launch time $T$ be zero or one, also each with probability one half. In model A, high quality launches immediately and zero quality launches later. In model B, the pairing is reversed. Both models have identical innovation counts and identical marginal laws of $Q$ and $T$. Nevertheless,
\[
 \E_A[\beta^TQ]=1,\qquad \E_B[\beta^TQ]=\beta,
\]
which differ whenever $\beta<1$. The relevant joint timing-quality information is missing. Replacing an expectation of their product by the product of marginal expectations would silently impose independence.

If only $0\leq G_k\leq\overline G_k$ is known while research and success probabilities are identified, a basic sharp value interval under otherwise unrestricted values is
\[
 -\sum_kc_kr_k\ \leq W(r)\leq
 \sum_k\overline G_kp_k(r_k)-\sum_kc_kr_k.
\]
Endpoints are attained by choosing all values at zero or their upper limits. Sharper policy comparisons require restrictions coupling the baseline and policy health values, not merely separate intervals optimized independently.

\section*{Remaining work}
The research technology and welfare weights are assumed known here. Real pricing reforms also alter access, insurance, prescribing, trial selection, quality and timing. Revenue variation can identify some private incentive responses under a credible design without identifying every health externality or the global optimal reward. A complete solution needs those joint causal responses, a feasible financing policy and uncertainty analysis over success and value. The exact allocation formula and counterexamples clarify what an empirical solution must measure, while leaving the full catalogue target unresolved.

\begin{thebibliography}{9}
\bibitem{al} D. Acemoglu and J. Linn (October 2003), Market Size in Innovation: Theory and Evidence from the Pharmaceutical Industry, NBER Working Paper 10038. Working-paper version checked; distinct from its 2004 journal publication.
\url{https://www.nber.org/system/files/working_papers/w10038/w10038.pdf}.
\end{thebibliography}
\end{document}
